\section{Exact cyclotomic identities and pole-cancelled spectral jets}
\label{ir:sec:identities}

\subsection{Evaluating an integral certificate}
For $x\in\Q/\Z$, write
\[
 \LL{s}{x}=\Li_s(e^{2\pi i x}).
\]
At $x=0$, this means $\zeta(s)$. For $x\ne0$, the spectral function
is entire. One way to see this, and to fix its continuation, is the
finite residue-class formula
\begin{equation}\label{ir:eq:fourier-hurwitz}
 \LL{s}{a/q}=q^{-s}\sum_{u=1}^{q}e^{2\pi i au/q}\zeta(s,u/q).
\end{equation}
For $a\not\equiv0\pmod q$, the residues of the simple poles at
$s=1$ sum to zero. The standard meromorphic continuation of Hurwitz
zeta has no other spectral poles \cite{dlmf-hurwitz}.

For $\Re s>1$, absolute convergence gives
\begin{equation}\label{ir:eq:analytic-distribution}
 \sum_{py=x}\LL{s}{y}=p^{1-s}\LL{s}{x}.
\end{equation}
Indeed the root sum cancels every series index not divisible by $p$;
for index $pk$ it contributes $p e^{2\pi i kx}$.
Continuation then proves the meromorphic identity for every complex
$s$. Thus an integral polynomial certificate becomes an exact analytic
identity under
\begin{equation}\label{ir:eq:polylog-weights}
 e_x\longmapsto\LL{s}{x},\qquad A_p\longmapsto p^{1-s},
\end{equation}
where real positive $p$ fixes the exponential unambiguously. No integer
relation search or floating-point acceptance threshold occurs in this
step.

\begin{proposition}[A removable endpoint coefficient]
\label{ir:prop:endpoint-zero}
For $y\ne0$, the coefficient $c_{y,0}(\mathbf A)$ in its integral
normal form vanishes at $A_p=1$ for all $p$. Hence its polylogarithmic
specialization cancels the simple pole of $\zeta(s)$ at $s=1$.
\end{proposition}
\begin{proof}
At ordinary weights, the functional taking $e_0$ to $1$ and all
other grid symbols to $0$ annihilates every distribution row. For
$x=0$, both the root sum and $e_x$ contain exactly one copy of
$e_0$; for $x\ne0$, neither does. Apply the functional to the
normal form of $e_y$. The second assertion follows because
$p^{1-s}=1$ at $s=1$ and the coefficient is analytic in $s$.
\end{proof}
The vanishing need only be first order. It is not legitimate to infer
one vanishing factor for each prime merely from the number of prime
variables in a raw-point normal form.

\subsection{A three-row level-twelve identity}
In the level-$12$ free module, define
\[
 F_{12}=e_{1/12}+e_{5/12}+(A_3+1)e_{3/4}
        -A_2A_3(A_2-1)e_0.
\]
The following is an equality before taking the quotient:
\begin{equation}\label{ir:eq:cert12}
 F_{12}=r_{3,1/4}+A_3r_{2,1/2}+A_2A_3r_{2,0}.
\end{equation}
For example, the first row contains
$e_{1/12}+e_{5/12}+e_{3/4}-A_3e_{1/4}$; the next two rows cancel
$e_{1/4}$ and then $e_{1/2}$. This verifies the certificate by direct
inspection.

\begin{corollary}\label{ir:cor:level12}
For every complex $s$, with the removable value used at $s=1$,
\begin{align}
 \LL{s}{1/12}+\LL{s}{5/12}
 +(3^{1-s}+1)\LL{s}{3/4}
 =6^{1-s}(2^{1-s}-1)\zeta(s).
 \label{ir:eq:level12}
\end{align}
At $s=1$, the right-hand side is $-\log2$.
\end{corollary}
\begin{proof}
Evaluate \eqref{ir:eq:cert12} using
\eqref{ir:eq:polylog-weights}. The value at one follows from
$2^{1-s}-1=-(s-1)\log2+O((s-1)^2)$.
\end{proof}

\subsection{A four-row level-fifteen identity}
Set
\begin{equation}\label{ir:eq:U-set}
 U=\{1/15,4/15,7/15,2/5,4/5\}.
\end{equation}
In the level-$15$ free module define
\begin{align}
 F_{15}={}&e_{8/15}-\sum_{u\in U}e_u
 -(A_3+1)e_{3/5}-(A_5+1)e_{2/3}
 -(1-A_3A_5)e_0.
 \label{ir:eq:F15}
\end{align}
A short exact certificate is
\begin{equation}\label{ir:eq:cert15}
 F_{15}=r_{3,3/5}-r_{5,1/3}-A_5r_{3,0}-r_{5,0}.
\end{equation}
The roots in the first two rows are respectively
\[
 \{1/5,8/15,13/15\},\qquad
 \{1/15,4/15,7/15,2/3,13/15\}.
\]
The $13/15$ terms cancel. The last two rows remove $1/3$ and
$1/5$, and leave precisely \eqref{ir:eq:F15}. In particular the
coefficient of $e_0$ in the normal form of $e_{8/15}$ is
$1-A_3A_5$, not $(A_3-1)(A_5-1)$.

\begin{theorem}[All-order level-fifteen identity]
\label{ir:thm:level15}
For every complex spectral order $s$,
\begin{align}
 \LL{s}{8/15}-\sum_{u\in U}\LL{s}{u}
 &-(3^{1-s}+1)\LL{s}{3/5}
  -(5^{1-s}+1)\LL{s}{2/3}\nonumber\\
 &=(1-15^{1-s})\zeta(s).
 \label{ir:eq:level15}
\end{align}
Both sides extend holomorphically through $s=1$, where the right-hand
side is $\log15$.
\end{theorem}
\begin{proof}
Apply \eqref{ir:eq:polylog-weights} to the four-row identity.
The endpoint value follows from
$1-15^{1-s}=(s-1)\log15+O((s-1)^2)$ and the residue one of
$\zeta(s)$.
\end{proof}
The theorem supplies infinitely many exact identities, including
ordinary positive integer weights, negative orders, and derivatives
in spectral order. It is generated by classical distribution rather
than by a new numerical relation conjecture.

A compatible level-$30$ normal form follows from the extra row
$r_{2,1/15}$:
\begin{align}
 e_{1/30}={}&(A_2-1)e_{1/15}-(A_3+1)e_{3/5}
 -(A_5+1)e_{2/3}\nonumber\\
 &-e_{4/15}-e_{7/15}-e_{2/5}-e_{4/5}
 +(A_3A_5-1)e_0.\label{ir:eq:level30-normal}
\end{align}
Its residual has the five-row certificate
\[
 r_{2,1/15}-r_{3,3/5}+r_{5,1/3}+A_5r_{3,0}+r_{5,0}.
\]
All rows can be read in the level-$30$ grid. This illustrates
compatibility of the straightener with an enlarged level and a
weight whose order-one specialization vanishes in a particular
coefficient.

\subsection{Every spectral derivative, with the pole term retained}
Let
\[
 \jj{j}{x}=\left.\frac{\partial^j}{\partial s^j}\LL{s}{x}
                   \right|_{s=1},\qquad x\ne0,
\]
and write the Laurent expansion
\begin{equation}\label{ir:eq:stieltjes}
 \zeta(1+t)=\frac1t+\sum_{m\ge0}\frac{(-1)^m\gamma_m}{m!}t^m.
\end{equation}
For a positive prime $p$ and $N\ge0$, define the explicit finite operator
\begin{equation}\label{ir:eq:T-operator}
 \mathcal T_{p,N}(x)=\jj{N}{x}
 +\sum_{j=0}^{N}\binom Nj(-\log p)^j\jj{N-j}{x}.
\end{equation}
It is the $N$th derivative of
$(p^{1-s}+1)\LL{s}{x}$ at $s=1$.

\begin{theorem}[A level-fifteen Stieltjes--polylogarithm jet family]
\label{ir:thm:jets15}
For every integer $N\ge0$, put $\ell=\log15$. Then
\begin{align}
 &\jj{N}{8/15}-\sum_{u\in U}\jj{N}{u}
 -\mathcal T_{3,N}(3/5)-\mathcal T_{5,N}(2/3)\nonumber\\
 &\qquad =(-1)^N\left\{\frac{\ell^{N+1}}{N+1}
 -\sum_{j=1}^{N}\binom Nj\ell^j\gamma_{N-j}\right\}.
 \label{ir:eq:jets15}
\end{align}
The sum on the right is empty at $N=0$.
\end{theorem}
\begin{proof}
All nonzero-root terms of \eqref{ir:eq:level15} are entire in $s$,
so differentiate the holomorphic continuation $N$ times. The weighted
terms give \eqref{ir:eq:T-operator} by Leibniz's rule. For the right
side use
\[
 1-e^{-\ell t}=\sum_{j\ge1}\frac{(-1)^{j+1}\ell^j}{j!}t^j
\]
and multiply by \eqref{ir:eq:stieltjes}. The pole term contributes
$(-1)^N\ell^{N+1}/(N+1)$ to the $N$th derivative. The regular
Laurent coefficients contribute
$(-1)^{N+1}\sum_{j=1}^N\binom Nj\ell^j\gamma_{N-j}$.
This is the stated formula.
\end{proof}

The first derivative gives the especially compact identity
\begin{align}
 \jj{1}{8/15}-\sum_{u\in U}\jj{1}{u}
 -2\jj{1}{3/5}-2\jj{1}{2/3}
 &+\log3\,\jj{0}{3/5}+\log5\,\jj{0}{2/3}\nonumber\\
 &=\gamma_0\log15-\tfrac12\log^2 15.
 \label{ir:eq:first-jet}
\end{align}
For $N=2$, the scalar right-hand side is
\[
 \frac{\log^3 15}{3}-\gamma_0\log^2 15-2\gamma_1\log15.
\]
These terms would be wrong if one substituted $s=1$ into the
coefficient of $\zeta(s)$ before cancelling its pole.

More generally, if an analytic coefficient $c(t)$ satisfies $c(0)=0$,
then
\begin{equation}\label{ir:eq:general-pole-jet}
 \left.\frac{d^N}{dt^N}\{c(t)\zeta(1+t)\}\right|_{t=0}
 =\frac{c^{(N+1)}(0)}{N+1}
 +\sum_{j=0}^{N}\binom Nj c^{(j)}(0)(-1)^{N-j}\gamma_{N-j}.
\end{equation}
This coefficient identity gives a general implementation rule for every
raw-point certificate, not just the displayed level.

\subsection{An elementary endpoint and branch check}
For $0<x<1$, the chosen branch gives
\[
 \jj{0}{x}=-\log(2\sin\pi x)+i\pi(1/2-x).
\]
Taking real parts of \eqref{ir:eq:level15} at $s=1$ gives
\[
 \frac{\prod_{u\in U}(2\sin\pi u)
 (2\sin(3\pi/5))^2(2\sin(2\pi/3))^2}
 {2\sin(8\pi/15)}=15.
\]
Since $\sin(8\pi/15)=\sin(7\pi/15)$, it simplifies to
\begin{equation}\label{ir:eq:sine-product}
 (2\sin(\pi/15))(2\sin(4\pi/15))
 (2\sin(2\pi/5))^3(2\sin(\pi/5))=5.
\end{equation}
All factors are positive, so exponentiating the real logarithmic identity
introduces no ambiguity. The imaginary parts cancel by the displayed
linear argument formula. This is an elementary specialization and branch
check, not a claim of novelty for the trigonometric product itself.

\subsection{A division-free Hurwitz--Stieltjes translation}
The same polynomial certificates evaluate under
\[
 A_p=p^s,\qquad e_x\mapsto Z_x(s),\qquad
 Z_x(s)=\begin{cases}\zeta(s,x),&0<x<1,\\\zeta(s,1),&x=0.\end{cases}
\]
The distribution identity here is
$\sum_{py=x}Z_y(s)=p^sZ_x(s)$; it follows by splitting the Hurwitz
series into residue classes, including the endpoint convention.
Let $c_b(t)$ be the normal-form coefficient of $e_b$ for $e_x$ after
$A_p=p^{1+t}$, and put $C(t)=\sum_b c_b(t)$.
Residue comparison gives $C(0)=1$. If
$\gamma_j(b^*)$ denotes the generalized Stieltjes coefficient with
$b^*=b$ for $b\ne0$ and $b^*=1$ for $b=0$, expansion of the exact
normal form gives, for all $N\ge0$,
\begin{align}
 \gamma_N(x^*)={}&
 \frac{(-1)^N C^{(N+1)}(0)}{N+1}
 +\sum_{b\in\B_q}\sum_{j=0}^{N}
 (-1)^j\binom Nj c_b^{(j)}(0)\gamma_{N-j}(b^*).
 \label{ir:eq:hurwitz-jets}
\end{align}
To verify the sign, multiply each Laurent series by $c_b(t)$, compare
$t^N$, and multiply by $(-1)^NN!$. The pole part is $C(t)/t$ and
contributes the first term. Thus even the underived constants require
a coefficient derivative. This is the raw-basis version of the
manuscript's pole-corrected jet philosophy, now with division-free
integral certificates for the coefficients.

Every coefficient derivative in \eqref{ir:eq:hurwitz-jets} is a finite
integer combination of products of logarithms of primes and integer
prime powers. There is no need to differentiate a solved matrix inverse
or to redo a relation search separately at each derivative order.
